% ============================================================
% Digit-Partitioning Primes and the Alignment Formula
% ============================================================

\documentclass[12pt]{amsart}

\usepackage{amsmath,amssymb,amsthm}
\usepackage{booktabs}
\usepackage{microtype}
\usepackage{url}

\newtheorem{theorem}{Theorem}
\newtheorem{lemma}[theorem]{Lemma}
\newtheorem{proposition}[theorem]{Proposition}
\newtheorem{corollary}[theorem]{Corollary}
\newtheorem{conjecture}[theorem]{Conjecture}
\theoremstyle{definition}
\newtheorem{definition}[theorem]{Definition}
\theoremstyle{remark}
\newtheorem{remark}[theorem]{Remark}

\title{Digit-Partitioning Primes and the Alignment Formula}

\author{Alexander S.\ Petty}

\email{alexander.petty@gmail.com}
\date{March 2020 (revised April 2026)}
\subjclass[2020]{11A63, 11B83}

\begin{document}
\begin{abstract}
For a prime $p$ not dividing a positional base $b$, the fractions
$\{k/p : 1 \le k \le p{-}1\}$ partition into cosets of the subgroup
$\langle b \rangle \le (\mathbb{Z}/p\mathbb{Z})^{*}$, each coset
determining a family of repetends related by cyclic permutation. We
say that $p$ is \emph{digit-partitioning} in base $b$ if no two
fractions from the full set $\{k/p\}$ share a digit at any common
repetend position unless they are identical.

We prove that $p$ is digit-partitioning in base $b$ if and only if
$p \le b + 1$. The proof rests on the injectivity of the digit
function $r \mapsto \lfloor br/p \rfloor$ on $\{1, \ldots, p{-}1\}$.

For every digit-partitioning prime, the repetend alignment of
$n = pm$ (with $m$ a $b$-smooth integer) equals exactly
$(2m{-}1)/(pm{-}1)$. This universal formula, combined with the
golden threshold analysis of the companion paper, confirms that the
golden ratio selects $p = 3$ uniquely among primes
$p \ge 3$: the alignment limit $2/p$ exceeds
$1/\varphi$ only when $p < 2\varphi \approx 3.236$,
which among primes means $p \in \{2, 3\}$.
\end{abstract}
\maketitle

% ============================================================
\section{Introduction}

This paper studies the digit function $\delta(r) = \lfloor br/p \rfloor$ and its structural invariants.

In a companion paper~\cite{paper1}, we proved that the repetend
alignment of $n = 3m$ (with $m$ a $10$-smooth integer) equals
$(2m{-}1)/(3m{-}1)$, and that the golden ratio is the unique
self-referential threshold classifying such integers. The formula
arose from the observation that non-matching fractions contribute
zero position-wise alignment with the base repetend.

This paper identifies the exact condition under which zero
cross-alignment holds for a general prime $p$ in a general base $b$.
The condition is $p \le b + 1$.
When this holds, the alignment formula $(2m{-}1)/(pm{-}1)$
generalizes without modification. Among primes $p \ge 3$,
the alignment limit $2/p$ exceeds $1/\varphi$ only for
$p = 3$.

% ============================================================
\section{Setting}

Let $b \ge 2$ be a positional base and $p$ a prime not dividing $b$.
The multiplicative order $L = \operatorname{ord}_{p}(b)$ is the
length of the repetend of $1/p$ in base $b$~\cite{hardy}. The group
$(\mathbb{Z}/p\mathbb{Z})^{*}$ has order $p - 1$, and the cyclic
subgroup $\langle b \rangle$ has order $L$, partitioning
$\{1, \ldots, p{-}1\}$ into $(p{-}1)/L$ cosets. Within each coset,
the repetends of $k/p$ are cyclic permutations of a common string.

\begin{definition}
An integer $m$ is \emph{$b$-smooth} if every prime factor of $m$
divides $b$.
\end{definition}

\begin{definition}
For $n = pm$ with $m$ a $b$-smooth integer, the \emph{repetend
alignment} $\alpha(n)$ is the mean, over $k \in \{1, \ldots, n{-}1\}$,
of the digit-match indicator $a(k/n,\,1/n)$ defined as follows:
if $k/n$ terminates, $a = 1$; if $k/n$ has a repetend of the
same length as $1/n$, $a$ is the fraction of positions (in the
standard eventually periodic expansion) at which the two repetends
share the same digit; otherwise $a = 0$.
\end{definition}

% ============================================================
\section{The Digit Function}

The central object is the digit function that maps remainders to
digits during long division. The floor function $\lfloor br/p \rfloor$
partitions the integers $\{1, \ldots, p{-}1\}$ into contiguous bins of
nearly equal size~\cite{fraenkel}.

\begin{definition}
For a prime $p$ not dividing $b$, the \emph{digit function} is
\[
\delta(r) = \left\lfloor \frac{br}{p} \right\rfloor,
\qquad r \in \{0, 1, \ldots, p{-}1\}.
\]
This is the digit produced when the remainder $r$ is processed in one
step of long division of any fraction with denominator $p$ in base $b$.
\end{definition}

\begin{lemma}\label{lem:injective}
If $p \le b + 1$, then $\delta$ is injective on $\{1, \ldots, p{-}1\}$.
\end{lemma}

\begin{proof}
We treat two cases.

\medskip\noindent
\emph{Case $p = b + 1$.}\quad
For $r \in \{1, \ldots, b\}$,
\[
\delta(r)
= \left\lfloor \frac{br}{b+1} \right\rfloor
= \left\lfloor r - \frac{r}{b+1} \right\rfloor
= r - 1,
\]
since $0 < r/(b{+}1) < 1$ for $1 \le r \le b$. The map
$r \mapsto r - 1$ is a bijection from $\{1, \ldots, b\}$ to
$\{0, \ldots, b{-}1\}$.

\medskip\noindent
\emph{Case $p \le b$.}\quad
For distinct $r, s \in \{1, \ldots, p{-}1\}$ with $r < s$, set
$D = b(s{-}r)/p$. From $s - r \ge 1$ and $p \le b$, we have
$D \ge b/p \ge 1$. We claim $D > 1$. Suppose to the contrary
that $D = 1$, so $b(s - r) = p$. Since $p$ is prime, $b \ge 2$,
and $s - r \ge 1$, this forces $s - r = 1$ and $b = p$,
contradicting $\gcd(b, p) = 1$. Therefore $D > 1$, so
$bs/p > br/p + 1$, and consequently
$\lfloor bs/p \rfloor \ge \lfloor br/p \rfloor + 1 > \lfloor br/p \rfloor$.
Thus $\delta(r) \ne \delta(s)$.
\end{proof}

\begin{remark}
When $p > b + 1$, the digit function maps $p - 1 > b$ distinct
remainders into $b$ digit values. By the pigeonhole
principle~\cite{stanley}, $\delta$ cannot be injective, and
collisions are unavoidable.
\end{remark}

% ============================================================
\section{The Digit-Partitioning Property}

\begin{definition}
A prime $p$ is \emph{digit-partitioning} in base $b$ if no two
distinct fractions in $\{k/p : 1 \le k \le p{-}1\}$ share a digit
at any common repetend position, unless their repetends are
identical.
\end{definition}

\begin{theorem}\label{thm:characterization}
Let $p$ be a prime not dividing $b$. Then $p$ is digit-partitioning
in base $b$ if and only if $p \le b + 1$.
\end{theorem}

\begin{proof}
At position $j$ of the repetend of $k/p$, the digit is
$\delta(b^{j}k \bmod p)$. For distinct fractions $k/p$
and $\ell/p$ with $k \not\equiv \ell \pmod{p}$, the
arguments $b^{j}k$ and $b^{j}\ell$ are distinct modulo $p$
(since $b^{j}$ is invertible).

If $p \le b + 1$, then $\delta$ is injective on
$\{1, \ldots, p{-}1\}$ by Lemma~\ref{lem:injective}, so
distinct arguments produce distinct digits at every
position. Therefore $p$ is digit-partitioning.

If $p > b + 1$, then $\delta$ maps $p - 1 > b$ remainders
into $b$ values. By the pigeonhole principle, there exist
distinct $r, s \in \{1, \ldots, p{-}1\}$ with
$\delta(r) = \delta(s)$. The first repetend digit of
$r/p$ equals the first repetend digit of $s/p$, so $p$ is
not digit-partitioning.
\end{proof}

% ============================================================
\section{The Universal Alignment Formula}

\begin{proposition}\label{lem:cycle-partition}
For $n = pm$ with $m$ a $b$-smooth integer and $p$ not dividing $b$,
the non-terminating fractions in $\{k/n : 1 \le k \le n{-}1\}$
partition into $(p{-}1)/L$ families of $mL$ fractions each,
where $L = \operatorname{ord}_{p}(b)$. Within each family, all
fractions share the same repetend up to cyclic shift. Fractions in
different families have distinct repetends.
\end{proposition}

\begin{proof}
The $m - 1$ multiples of $p$ in $\{1, \ldots, pm{-}1\}$ give
terminating fractions. The remaining $(p{-}1)m$ values of $k$ are
not divisible by $p$.

Choose $t$ with $m \mid b^{t}$, and write $b^{t} = mu$
where $\gcd(u, p) = 1$ (possible since $m$ is $b$-smooth
and $p \nmid b$). Then
\[
\frac{k}{pm} = \frac{ku}{p \cdot b^{t}},
\]
so after $t$ base-$b$ digits, the repetend of $k/(pm)$
is exactly the repetend of $(ku \bmod p)/p$.
Since multiplication by $u$ permutes
$(\mathbb{Z}/p\mathbb{Z})^{*}$, the nonterminating
repetends of $\{k/(pm)\}$ are in bijection with
the repetends of $\{r/p : 1 \le r \le p{-}1\}$.

The residues $1, \ldots, p{-}1$ partition into
$(p{-}1)/L$ cosets of $\langle b \rangle$ in
$(\mathbb{Z}/p\mathbb{Z})^{*}$. Residues in the
same coset produce cyclic shifts of a single
repetend. Residues in different cosets produce
distinct repetends (since their orbits under
multiplication by $b$ are disjoint).

Each residue class modulo $p$ contains exactly $m$ elements of
$\{1, \ldots, pm{-}1\} \setminus p\mathbb{Z}$, so each coset
accounts for $mL$ non-terminating fractions.
\end{proof}

\begin{theorem}\label{thm:alignment}
If $p \le b + 1$ and $m \ge 1$ is $b$-smooth, then
\[
\alpha(pm) = \frac{2m - 1}{pm - 1}.
\]
\end{theorem}

\begin{proof}
The $pm - 1$ fractions partition into three groups:

\begin{enumerate}
\item[\textup{(A)}]
$m - 1$ terminating fractions (multiples of $p$). Each contributes
alignment~$1$.

\item[\textup{(B)}]
$m$ fractions whose repetend matches that of $1/(pm)$.
Since $1/p$ lies in a single coset family,
Proposition~\ref{lem:cycle-partition} gives exactly $m$
fractions in $\{k/(pm)\}$ with the same repetend.
Each contributes alignment~$1$.

\item[\textup{(C)}]
$(p{-}2)m$ remaining non-terminating fractions.
Since $p$ is digit-partitioning, the digit function is
injective (Theorem~\ref{thm:characterization}), so
fractions from different cosets share no digit at
any common position. Each contributes alignment~$0$.
\end{enumerate}

Total: $(m{-}1) + m = 2m - 1$ aligned out of $pm - 1$.
\end{proof}

\begin{corollary}\label{cor:limit}
For fixed $p$, the alignment $\alpha(pm) \to 2/p$ as
$m \to \infty$.
\end{corollary}

% ============================================================
\section{Enumeration}

For each base $b$, the digit-partitioning primes are precisely the
primes $p \le b + 1$ that do not divide $b$.

\begin{table}[h]
\centering
\begin{tabular}{rll}
\toprule
Base $b$ & Digit-part.\ primes ($p \le b{+}1$, $p \nmid b$) & Alignment limit $2/p$ \\
\midrule
$4$  & $3, 5$               & $0.667,\; 0.400$ \\
$6$  & $5, 7$               & $0.400,\; 0.286$ \\
$8$  & $3, 5, 7$            & $0.667,\; 0.400,\; 0.286$ \\
$10$ & $3, 7, 11$           & $0.667,\; 0.286,\; 0.182$ \\
$12$ & $5, 7, 11, 13$       & $0.400,\; 0.286,\; 0.182,\; 0.154$ \\
$16$ & $3, 5, 7, 11, 13, 17$& $0.667,\; \ldots,\; 0.118$ \\
\bottomrule
\end{tabular}
\caption{Digit-partitioning primes for selected bases and their
alignment limits $2/p$.}
\label{tab:bases}
\end{table}

In base~$10$, the digit-partitioning primes are $3$, $7$, and $11$.
These correspond to the three classical mechanisms identified in
earlier analysis:
\begin{itemize}
\item $p = 3$: $\operatorname{ord}_{3}(10) = 1$ (single-digit
      repetends).
\item $p = 7$: $\operatorname{ord}_{7}(10) = 6 = p{-}1$ (cyclic
      number, primitive root).
\item $p = 11$: $\operatorname{ord}_{11}(10) = 2$, with
      $10 \equiv -1 \pmod{11}$ (nines-complement partition).
\end{itemize}

These mechanisms are unified by the single condition
$p \le b + 1$.

% ============================================================
\section{Connection to the Golden Ratio}

The alignment limit $2/p$ exceeds $1/\varphi$ only when
$p < 2\varphi \approx 3.236$. Among primes $p \ge 3$,
only $p = 3$ meets this condition. (For $p = 2$, the
alignment limit is $1$, trivially above any threshold.)

The self-referential characterization developed
in~\cite{paper1} (Sections~5 and~6) applies: the
golden ratio's minimal polynomial
$\tau^{2} + \tau - 1$ divides the self-referential cubic
if and only if $p = 3$, producing the fourth-power
identity $\tau^{4} = 2 - 3\tau$. The algebra is proved
in the companion paper; this section uses the result.

The present paper strengthens this result. The formula
$(2m{-}1)/(pm{-}1)$ is not a special case tied to base~$10$ or to any
particular mechanism. It holds universally for all digit-partitioning
primes in all bases, and the characterization $p \le b + 1$ shows
that the class of such primes is determined by a single inequality.
Among primes $p \ge 3$ in this class, $p = 3$ is the
unique value with alignment limit above $1/\varphi$.

% ============================================================
\section{Remarks}

\subsection*{Structural manifestations}

Although the characterization is a single inequality, the
digit-partitioning property manifests through different structural
mechanisms depending on the relationship between $b$ and $p$:

\begin{itemize}
\item When $b \equiv 1 \pmod{p}$: repetends have length one, and
      each fraction uses a unique single digit.
\item When $p = b{+}1$ (so $b \equiv -1 \pmod{p}$): repetends
      have length two, and cosets correspond to $(b{-}1)$-complement
      pairs, partitioning the digit set completely.
\item When $b$ is a primitive root modulo $p$ and $p \le b{+}1$:
      there is a single coset, and the repetend is a cyclic number
      with all distinct digits.
\end{itemize}

These are consequences of the injectivity of $\delta$, not
independent conditions.

\subsection*{The boundary case $p = b + 1$}

The case $p = b + 1$ is the tightest: $\delta$ maps $b$ remainders
into $b$ digit values, and injectivity holds by a narrow margin. In
base~$10$, this gives $p = 11$, whose five nines-complement pairs
partition all ten digits exactly. The boundary case
always produces the complement-pair structure, since
$b + 1 = p$ implies $b \equiv -1 \pmod{p}$.

\subsection*{Open questions}

For $p > b + 1$, the alignment formula acquires correction terms from
digit collisions. The corrections are governed by the squared bin sizes
of the digit function and have a closed form for primitive-root primes.

% ============================================================
\begin{thebibliography}{9}

\bibitem{paper1}
A.~S.~Petty,
Why the golden ratio selects the prime three,
research note, 2020.

\bibitem{fraenkel}
A.~S.~Fraenkel,
The bracket function and complementary sets of integers,
\emph{Canad.\ J.\ Math.} \textbf{21} (1969), 6--27.

\bibitem{stanley}
R.~P.~Stanley,
\emph{Enumerative Combinatorics}, vol.~1,
2nd ed., Cambridge University Press, 2012.

\bibitem{hardy}
G.~H.~Hardy and E.~M.~Wright,
\emph{An Introduction to the Theory of Numbers},
6th ed., Oxford University Press, 2008.

\end{thebibliography}

\end{document}
